\subsection{Evaluating the compatible radical by ring transitions.}
The context tests of Theorem~\ref{thm:upper-levitzki-quotient} can be
replaced by transitions between explicit ordinary rings. This evaluates
the same coefficient ideal; it does not classify the remaining projectives.

\begin{proposition}\label{prop:upper-transition-radical}
For $e\in U(Q)$, the compatible coefficient ideal of
Theorem~\ref{thm:upper-levitzki-quotient} satisfies
\[
 H_e=\bigcap_{\substack{f\in U(Q)\\e\leq f}}
       t_{e,f}^{-1}\bigl(\operatorname{Lev}(F(f))\bigr).
\]
For any $q\in Q$, put
\[
 E(q)=\begin{pmatrix}1&q\\0&1\end{pmatrix},
 \qquad
 j_q:F(q)\longrightarrow A_{E(q)}(T),
\]
where $j_q(x)$ has coefficient $x$ in position $(1,2)$ and zero
amplitude in the other specified groups. Then
\[
 H_q=\bigcap_{\substack{D\in M_2(Q)\\D^2=D,\ E(q)\leq D}}
 (t_{E(q),D}j_q)^{-1}\bigl(\operatorname{Lev}(A_D(T))\bigr).
\]
All target rings are unital. If $Q$ is finite, the second intersection
has at most $|Q|^4$ terms, without any finiteness assumption on $F$.
Computing the ordinary radicals remains a separate problem.
\end{proposition}

\begin{proof}
We first recall the ordinary transfer argument underlying the corner
result of Rich~\cite[Corollary~1, p.~57]{Rich}. In a Morita context with
diagonal rings $A,B$ and off-diagonal bimodules $M,N$, let $I$ be a
locally nilpotent ideal of $B$. The additive group $MIN$ generated by
$mzn$, with $z\in I$, is an ideal of $A$. Expand a finite subset using
finitely many terms $m_i z_i n_i$. Products of these terms have the form
\[
 m_{i_1}z_{i_1}(n_{i_1}m_{i_2})z_{i_2}\cdots
 (n_{i_{s-1}}m_{i_s})z_{i_s}n_{i_s}.
\]
The finitely many elements $z_i(n_i m_j)$ belong to $I$ and generate
a nilpotent subring. Thus sufficiently long products vanish, proving
$M\operatorname{Lev}(B)N\subseteq\operatorname{Lev}(A)$.

Necessity of the first formula follows from transition stability of
$H$ and $H_f\subseteq\operatorname{Lev}(F(f))$. For sufficiency, suppose
$x\in F(e)$ passes every stated upward test. Consider any defining
context $u\in F(p)$, $v\in F(r)$ with $per\leq h\in U(Q)$. Set
\[
 a=hpe,\qquad b=erh,\qquad f=e\vee ba=e\vee erhpe.
\]
We have $ab\leq h$, $e(ba)=ba=(ba)e$, and
$(ba)^2=b(ab)a\leq bha=ba$. Consequently $f^2=f$ and $f\geq e$.
Moreover
\[
 ha=a=af,\qquad fb=b=bh.
\]
It follows that
\[
 D=\begin{pmatrix}h&a\\b&f\end{pmatrix}
\]
is an upper multiplicative idempotent. Its ordinary fiber ring
$A_D(T)$ has diagonal corners $F(h),F(f)$ and off-diagonal groups
$F(a),F(b)$. Transport $u,v$ to these off-diagonal groups; this is
possible because $p\leq a$ and $r\leq b$. By hypothesis
$t_{e,f}(x)\in\operatorname{Lev}(F(f))$. The ordinary transfer argument
therefore places their threefold product in
$\operatorname{Lev}(F(h))$. Multiplicative naturality identifies it
with $t_{per,h}(uxv)$: the support identities $ae=a=af$ and
$eb=b=fb$ justify the transports without changing their order.
Thus every original context test passes. No common upper bound for
all upper idempotents has been assumed.

Finally $E(q)^2=E(q)\geq I_2$. Finite-matrix compatibility in
Theorem~\ref{thm:upper-levitzki-quotient} gives
\[
 x\in H(T)_q\quad\Longleftrightarrow\quad
 j_q(x)\in H(M_2(T))_{E(q)}.
\]
Apply the first formula to $M_2(T)$ and $E(q)$. Its upper fiber rings
and transitions are precisely $A_D(T)$ and $t_{E(q),D}$, giving the
second formula. There are at most $|Q|^4$ support matrices when $Q$
is finite. The group $F(0)$ is retained throughout; the lower-left
entry of $j_q(x)$ is its zero element, not an assertion that the
group is zero.
\end{proof}
